\section{Methods}

\subsection{Task, and why the unit of structure is a complex}

The model receives a hit molecule $x$ and a structured description $c$ of that molecule's
environment in its target, and emits an intent sentence $z$ and an analogue $y$. Critically, $c$ is
\emph{ligand-specific}: one block per ligand heavy atom, sixteen fields giving the nearest protein
atom and its distance, a typed interaction with hydrogen-bond angle, a second polar partner, water
counts and nearest water, a lining description, an open fraction and a burial reading (full field
list in the SI). A target-level pocket is therefore \emph{not} sufficient input: knowing the protein
buys nothing for a ligand whose pose we have never computed. The unit of work is the pair (molecule, target) --- a
complex --- and the cost of the structural conditioning is one structure per molecule, not one per
target. This single fact drives the dataset economics reported below.

Where a crystal structure of the complex exists we use it. Where it does not --- which is the case
for essentially every molecule in the literature corpus --- we cofold the complex with
Boltz~\citep{wohlwend2024boltz1,passaro2025boltz2} at one diffusion sample and derive the table from
the returned structure. This bounds what may be claimed from it, and we state the bound rather than
hoping a referee misses it:

\begin{quote}
Where no crystal structure exists, the complex is a single-sample cofold. Its contact table is a
hypothesis about which residues line the ligand, not a measurement of distances. No distance taken
from a cofolded table is quoted as evidence in this paper, and covalent geometry is never read from
a cofold, because an unconstrained cofold forms no bonds and a warhead lands wherever the sampler
puts it.
\end{quote}

\paragraph{Covalent typing.} A covalent bond is asserted only where the table types the row
\texttt{COVALENT BOND}, requiring the modelled heavy-atom separation to lie within 0.35\,\AA\ of the
ideal length for that element pair (1.81\,\AA\ for C--S; the SARS-CoV-2 main protease complex
7JQ2 measures exactly 1.81\,\AA\ at CYS145 and is our calibration case). Separations of
2.16--2.71\,\AA\ are typed and disclosed as \emph{probable} and are never cited as a bond; anything
beyond 2.71\,\AA\ is a non-bonded contact however a sentence describes it. This matters because an
audit of an earlier generation of our own labels found that \m{26\textrm{--}40\%} of clauses calling
a residue a ``covalent anchor'' were in fact citing a 2.7--3.8\,\AA\ non-bonded contact --- a
failure of the text layer, not of the structure layer, which had typed those rows correctly all
along. Every clause matching \texttt{covalent|bond|anchor|warhead|adduct} is now gated on a
\texttt{COVALENT BOND} row in that pair's own table, and the gate runs irrespective of whether the
sentence quotes a distance.

\subsection{Corpora}

\begin{table}[H]\centering\small
\begin{tabular}{llll}
\toprule
corpus & what it is & scale & intent written by \\
\midrule
PocketPair-9k/10k/14k & PDB-derived analogue pairs with pockets & 7{,}692 / 9{,}453 / 12{,}392 rows & an agent, from the pair \\
PocketPair-21k & ChEMBL within-assay pairs, cofolded, with $\Delta$potency & 3{,}992 complexes & an agent, from the pair \\
PocketPair-Literature & design intents quoted from paper text & 1{,}622 papers, 17{,}601 pairs & the paper's authors \\
\bottomrule
\end{tabular}
\caption{The three supervision regimes. The distinction in the last column is the one that matters
for interpreting our strongest result, and we return to it in \S\ref{sec:leakage}.}
\end{table}

\paragraph{Literature extraction.} For each paper we extract design \emph{episodes}: a reference
compound, one or more analogues, what changed, the stated reason, a verbatim evidence quote, the
endpoint in the paper's own words, and the direction of the outcome. Molecules are not recovered by
substructure search --- which finds the authors' actual pair only \m{32\%} of the time --- but by
ChEMBL's own curation under the paper's compound labels. A potency $\Delta$ is computed only between
two molecules measured in the \emph{same assay} of that paper. Under the final brief (v5.3), on all
\m{1{,}622} papers, the endpoint is stated in \m{100\%} of records, the reference source in
\m{100\%}, a rationale in \m{100\%} and an evidence quote in \m{100\%}; extraction-meta leakage is
\m{0.00\%}; repeated boilerplate in \texttt{what\_changed} is \m{0.0\%}; and edits spanning multiple
sites fell from \m{44.8\%} to \m{5.5\%} between brief versions.

\paragraph{Corpus defects, stated here rather than in a footnote.} These are SAR pairs, not matched
molecular pairs: the gate is a similarity threshold plus a same-assay potency delta, and the
extracted transformation field is descriptive only. A 100-paper QA round found molecules contradicted
by the text in \m{3.3\%} of records and reasons honest in \m{90\%} --- but \m{60.7\%} of those
reasons \emph{restate an outcome} rather than state an instruction, so only \m{39.3\%} of the corpus
is the supervision we actually want. Training on outcome restatements teaches post-hoc narration,
which is precisely the failure an independent cohort round measured at explanation-sound \m{0.08}.
We therefore report the instruction subset separately \TBD{13}{instruction-subset training};
hypothesis: the instruction subset trains a more useful $R$ at equal token count. Separately,
explained edits carry \emph{small} potency deltas, so selecting pairs by $|\Delta|$ systematically
discards them --- intent mining has to precede pair selection, not follow it.

\paragraph{Complex coverage.} The \m{17{,}601} literature pairs require \m{6{,}590} distinct
complexes across \m{658} targets, of which \m{27} were already held. \m{552} targets have a
sequence in ChEMBL and \m{106} do not, placing \m{1{,}995} complexes and \m{4{,}967} pairs out of
reach without another sequence source. \m{369} complexes exceed our \m{1{,}187}-residue cofold cap
and are recorded in a separate file rather than dropped --- an earlier run lost an entire major
target family to exactly this cap, silently. The remaining \m{4{,}199} complexes are queued in
descending order of pairs unlocked, so a partial run is the most valuable slice rather than a random
one: the first \m{1{,}000} complexes unlock \m{43.2\%} of the needed pairs and the first \m{3{,}000}
unlock \m{58.6\%}.

\subsection{Models: generated reasoning, extracted reasoning, and direct training}

Three training regimes are compared, and the paper's framing depends on keeping them distinct.

\emph{Direct LLM training} is the baseline: one adapter, one sequence, rationale and SMILES emitted
together, with nothing in the objective forcing the molecule to depend on the rationale and no handle
to intervene on under a mismatched-intent test. \emph{Generated reasoning} is our $R$, a separate
adapter emitting the intent from the complex alone and supplying $M$ at inference.
\emph{Extracted reasoning} is $R$'s supervision drawn from the literature rather than from the
agent-written intents of the crystal and potency corpora --- a difference that is not cosmetic, since
an agent writing an intent for a pair has the product molecule in view, whereas a paper's author
wrote the hypothesis without knowing what the reader would later be told. We treat agent-written intents as a \emph{weak upper bound} on supervision
quality and literature intents as the target, and the comparison between them is a result
(\S\ref{sec:agentic}), not a design choice.

\paragraph{Bases and finetuning.} The primary base is ChemDFM-v2.0-14B~\citep{zhao2024chemdfm};
ether0-24B~\citep{narayanan2025ether0}, Llama-3.1-8B-Instruct~\citep{dubey2024llama3},
SmileyLlama-8B~\citep{cavanagh2024smileyllama} and DESRO-8B~\citep{liu2026desro} are comparators.
Because these differ in scale (8B to 24B), cross-base rows are confounded with capacity and the only
clean contrast is each base against \emph{itself}, with and without our adapter. All arms use LoRA
($r{=}16$, $\alpha{=}32$, dropout 0.05, all seven projections; lr $10^{-4}$ cosine, batch 1 $\times$
accumulation 8, bf16) with one shared configuration across every cell of an ablation, and a single
fixed instruction format for every model rather than per-model chat templates, so that template is
never confounded with checkpoint; the blind arm substitutes ``No pocket structure is provided.'' for
the table. Loss is masked to the response --- the prompt runs a median of roughly 960 tokens against
roughly 100 for the answer --- and the length limit acts as a refusal threshold rather than a
truncation, because truncating would delete contact-table rows from precisely the arm the pocket
ablation is meant to favour.

\paragraph{Leakage controls.} For $M$, any SMILES line inside the provided reason is stripped and
the removal is asserted per row; the analogue SMILES, the analogue's PDB identifier and
analogue-only residue tokens are asserted absent from every prompt. Special-token identifiers are
pinned per vocabulary and re-verified at load.

\paragraph{Checkpoint selection, and the resulting caveat.} Training runs to at most 10 epochs with
early stopping on validation loss, patience 2, and the best checkpoint retained. Consequently, and
stated verbatim in every validation caption:

\begin{quote}
Because both the stopping epoch and the retained checkpoint were selected on the validation split,
every validation number reported here is a model-selection score, not a held-out test score, and is
optimistically biased by an unknown amount.
\end{quote}

\subsection{Splits}

Our held-out claims rest on a \textbf{molecule-level} holdout: neither the input nor the reference
analogue appears anywhere in training. On the 21k corpus this leaves \m{1{,}469} rows
(\m{553} on targets never seen in training, \m{916} on trained targets with new molecules), from
which \m{400} prompts per arm are drawn so that all arms answer identical inputs. Unseen-\emph{target}
holdout is infeasible here: target space repeats heavily across papers, so demanding unseen targets
leaves roughly \m{21} of them and the variance would be dominated by which handful of proteins
landed in the holdout. \TBD{14}{split reconciliation}: the corpus builder inherits a target-level
val fraction (32 val targets) while the holdout above is molecule-level; Methods must state one
authoritative sentence for the split the model was \emph{trained} under and the holdout it was
\emph{scored} on. We do not use random pair splits at all: sibling pairs off one complex share the
pocket, the numbering and a near-identical table, so a random split places them on both sides and
measures memorisation --- in a sibling project on the same pair machinery, \m{71\%} of a random
split's test pairs were exact duplicates of training pairs.

\subsection{Evaluation}

\paragraph{Molecular metrics.} \emph{Validity} is a SMILES extracted and parsed, over all rows, so
a model emitting no molecule scores 0 rather than NaN; it is a gate, not a result.
\emph{Uniqueness} is distinct canonical SMILES over valid ones, and functions here as a
collapse detector --- note that in our arms it \emph{anti}-correlates with quality, the mismatched
arm having both the highest uniqueness (\m{0.937}) and the worst molecules. \emph{TC$\rightarrow$input}
(ECFP4, 2048 bits) is a window check, not a quality score: reference analogues sit near \m{0.68},
which is a band and not a maximum, and \m{1.000} means the input was echoed.
\emph{TC$\rightarrow$reference} is the only metric whose ground truth is a real human design
decision, and it is necessary but not sufficient --- it penalises a different \emph{valid} analogue
--- so it is used only in paired contrasts between arms answering the same input.
\emph{Heavy-atom $\Delta$ MAE} catches the right-size/wrong-edit failure that Tanimoto hides.

\paragraph{The groundedness ladder.} ``Groundedness'' as commonly reported means only that a cited
residue exists in that pair's own table. That is rung L0 of a ladder, and the paper reports the
ladder rather than the headline: L0 existence; L1 residue and distance co-occurring in one row;
L2 interaction type matching the row's typing; L3 the covalent gate above; L4 \emph{site adjacency},
the cited residue contacting at least one atom of the changed set; L5 directional ownership, claims
about what is given up mapping to changed atoms; L6 \emph{atom ownership}, the row's atom name
mapped through the chemical-component graph to a specific input index inside the changed set. No
claim in this paper headlines a rung below L4. Each rung is reported as a pair --- pass rate
\emph{and} measurable fraction --- because the mapping fails on some pairs, and an earlier panel
scoring 6/6 against a rigorous 2/6 was partly a coverage gap rather than a disagreement.
Stereochemistry-only pairs are unscorable at L4 and are reported as unscored, never as failed.
\emph{Invented chemistry} --- a sentence asserting a contact absent from that pair's table --- is the
one rationale metric that needs no panel, and it is monotone across our arms.

\paragraph{Residue specificity, not selectivity.} The overlap between the residues a sentence names
and those the reference sentence names is reported as \emph{residue specificity}. We previously
called it selectivity, which invites the reading ``off-target selectivity'' and the demand for a
kinase panel; it is nothing of the kind. It is printed beside its zero-parameter baseline --- name
the nearest table residue --- because an arm at that baseline is no better than a distance sort.

\paragraph{Panels, and what they cannot do.} Sentence quality is scored both by model labellers
under a fixed rubric and by chemists. Panel-to-panel drift exceeds the between-arm differences we
measure, so every panel number is reported with its seeded controls read first, no comparison is
ever made across panels, and panel figures are never pooled with programmatic ones. We report
\emph{pooled multi-target $R^2$ nowhere}: between-target potency offsets inflate it even when
within-target ranking is random. Per-target averaged statistics with $n$ per target only.
